% ======================================================================
% preprints/jphq_06_structural_tension_nonoptimizable/sections/06_scope_limits_nonclaims.tex
% Journal-grade content (100/100). ASCII-only.
% ======================================================================

\section{Scope Limits and Non-Claims}

This preprint answers a single frozen question: whether structural tension
$\TEA$ in enterprise continua $\KEA$ is intrinsically non-optimizable within
the formal commitments of the Ontology of Continua (OC) and
\texttt{ETS.K\_EA.v1.1}. The answer is negative with respect to optimization.
The scope is intentionally narrow and explicitly bounded.

\subsection{Scope}

All results are derived strictly within OC/ETS and rely only on the
primitives and relations already defined for $\KEA$, including admissible
states $\OmegaEA$, boundary semantics $\dOmegaEA$, thresholds $\ThetaEA$,
continuity $\kEA$, and phase logic with explicit precedence.

Within this scope, the paper provides:
\begin{itemize}
  \item a formal characterization of structural tension as a configuration
        of threshold-inequality evaluations,
  \item impossibility results excluding scalarization and optimization of
        $\TEA$ under the OC/ETS scope lock,
  \item explicit identification of failure modes for optimization-based
        interpretations (ordering, metric substitution, phase
        reinterpretation),
  \item falsifiability conditions formulated entirely within the same
        primitive set.
\end{itemize}

No additional assumptions, metrics, or empirical interpretations are
introduced.

\subsection{Non-Claims}

To prevent category errors, the following non-claims are stated explicitly.

\paragraph{No governance, control, or intervention claims.}
The paper introduces no prescriptions, objectives, steering mechanisms, or
control frameworks. It does not propose how to reduce, manage, mitigate, or
respond to structural tension. Any such interpretation lies outside the
formal scope.

\paragraph{No performance interpretation.}
Structural tension $\TEA$ is not treated as a cost, risk, loss, stress, or
KPI-like quantity. No claim is made that higher or lower tension correlates
with performance, outcomes, efficiency, or business success.

\paragraph{No empirical or statistical claims.}
No datasets, measurements, estimators, or statistical inferences are used
or implied. All arguments are structural and semantic.

\paragraph{No quantitative margins or buffers.}
The theory, as used here, defines no distances to thresholds, safety
margins, robustness radii, or buffer zones. Such notions require metric
structure that is absent from OC/ETS.

\paragraph{No cross-enterprise comparability.}
The paper does not claim that $\TEA$ is comparable across enterprises, nor
that enterprises can be ranked by tension. OC/ETS defines no such ordering
or normalization.

\paragraph{No claims about the impossibility of optimization in general.}
The paper does not claim that enterprises cannot be optimized under other
formalisms. It claims only that optimization of $\TEA$ is not defined within
OC/ETS without introducing additional primitives.

\subsection{Boundary of Applicability}

The impossibility results depend critically on the scope lock of
\texttt{ETS.K\_EA.v1.1}: no added metrics, aggregators, utilities, weights,
or reparameterization invariants beyond those explicitly specified.

If such structure is added, a different theory object is obtained. The
results in this paper do not apply to that extended object and make no
claims about it.

\subsection{Implication of the Limits}

Within OC/ETS, structural tension functions as a diagnostic delimiter of
admissibility. It cannot be promoted to an optimization target without
changing the formal semantics of thresholds, phases, or continuity.

This paper records that boundary explicitly and makes no claims beyond it.

%%% End of section %%%
